\selectlanguage{lithuanian}
\sectionnonum{Santrauka}
\singlespacing
\listoffigures % TODO: dabar dar nera jokiu
\listoftables
\tableofcontents
\onehalfspacing
\sectionnonum{Santrumpos}

\small
\begin{tabular}{lp{0.72\textwidth}}
    ABI & \textit{Application Binary Interface} -- programų dvejetainė sąsaja \\
    ASID & \textit{Address Space Identifier} -- adresų erdvės identifikatorius \\
    CLINT & \textit{Core Local Interrupt Controller} -- branduolio vietinis pertraukčių valdiklis \\
    CSR & \textit{Control and Status Register} -- valdymo ir būsenos registras \\
    DMA & \textit{Direct Memory Access} -- tiesioginė prieiga prie atminties \\
    DTB & \textit{Device Tree Blob} -- dvejetainis įrenginių medžio aprašas \\
    DWARF & \textit{Debugging With Attributed Record Formats} -- derinimo informacijos formatas \\
    ELF & \textit{Executable and Linkable Format} -- vykdomųjų ir susiejimo failų formatas \\
    ISA & \textit{Instruction Set Architecture} -- instrukcijų rinkinio architektūra \\
    MMIO & \textit{Memory-Mapped Input/Output} -- į atminties adresų erdvę susieta įvestis ir išvestis \\
    PLIC & \textit{Platform-Level Interrupt Controller} -- platformos lygmens pertraukčių valdiklis \\
    PPN & \textit{Physical Page Number} -- fizinio puslapio numeris \\
    PTE & \textit{Page Table Entry} -- puslapių lentelės įrašas \\
    QPU & \textit{Quantum Processing Unit} -- kvantinių skaičiavimų įrenginys \\
    RAM & \textit{Random Access Memory} -- laisvosios prieigos atmintis \\
    SBI & \textit{Supervisor Binary Interface} -- supervisor režimo dvejetainė sąsaja \\
    UART & \textit{Universal Asynchronous Receiver/Transmitter} -- universalus asinchroninis imtuvas ir siųstuvas \\
\end{tabular}

\sectionnonum{Įvadas}

\textbf{TODO} Darbo tikslai:

\begin{enumerate}
    \item Išanalizuoti RISC-V architektūros, privilegijuoto vykdymo, atminties valdymo ir Linux paleidimo reikalavimus, reikalingus pilnai RV64 sistemos emuliacijai.

    \item Suprojektuoti ir realizuoti RV64 emuliatorių, gebantį vykdyti bazines, privilegijuotas ir atomines RISC-V instrukcijas bei emuliuoti Linux paleidimui būtinus įrenginius.

    \item Įgyvendinti eksperimentinį kvantinių skaičiavimų instrukcijų plėtinį ir paruošti programinės įrangos grandinę, leidžiančią šias instrukcijas naudoti vartotojo erdvės programose.

    \item Patikrinti emuliatoriaus teisingumą naudojant testų rinkinius ir įvertinti našumą naudojant eksperimentinį instrukcijų rinkinį.
\end{enumerate}
